\section{Experiments}
\label{sec:exp}

We design experiments to answer four questions: (Q1) Does \StereoSplat
recover metric scale while preserving photometric quality? (Q2) How does
each component contribute? (Q3) Do the display operators reduce temporal
artifacts? (Q4) What does the reconstruction look like qualitatively?

\paragraph{Datasets.}
We evaluate on three stereo datasets with diverse scene types.
\textbf{TartanAir}~\cite{wang2020tartanair} is a synthetic dataset with
ground-truth metric depth and pose across diverse environments, providing
an oracle upper bound. \textbf{DSEC}~\cite{gehrig2021dsec} is a driving
stereo dataset with calibrated rectified pairs and LiDAR depth
supervision. \textbf{KITTI Stereo 2015}~\cite{menze2015kitti} provides
sparse ground-truth disparities for evaluation. For all datasets we
precompute IGEV-Stereo disparities with
\texttt{preprocess\_disparity\_stereo.py} and convert to metric depth via
Eq.~\ref{eq:triang}.

\paragraph{Baselines.}
We compare against: (i) \StreamSplat~\cite{wu2026streamsplat} (monocular,
orthographic, no metric scale); (ii) \StereoSplat \emph{without} metric
depth (uses relative depth, keeps perspective rasterizer; isolates the
contribution of the stereo depth source); (iii) an oracle variant that
uses ground-truth metric depth (upper bound). All methods share the same
training schedule and architecture hyperparameters.

\paragraph{Metrics.}
We report PSNR, SSIM~\cite{wang2004ssim}, and LPIPS~\cite{zhang2018lpips}
on novel-view synthesis (left-image rendering and, where applicable,
right-image rendering). For metric scale we report the metric depth L1
error in meters, computed only on pixels with valid ground-truth depth.
For temporal consistency we report the per-pixel temporal standard
deviation across the sequence in static regions, following the
spatiotemporal-image protocol of~\cite{yun2025compensating}.

\subsection{Main Comparison (Q1)}

Table~\ref{tab:main} reports the main comparison. \StereoSplat achieves
metric depth L1 error of \TODO{X}\,m on TartanAir, \TODO{X}\,m on DSEC,
and \TODO{X}\,m on KITTI, while matching or surpassing \StreamSplat on
PSNR/SSIM/LPIPS. The no-metric-depth variant matches \StreamSplat
photometrically but cannot report a metric-depth error, confirming that
metric scale comes specifically from the stereo front end. The oracle
variant upper-bounds the metric quality and shows that the gap to perfect
depth is small.

\begin{table}[t]
\centering
\small
\caption{Main comparison. PSNR/SSIM/LPIPS on left-view rendering; depth
L1 in meters (lower is better). \TODO{numbers to be filled after
training runs complete.}}
\label{tab:main}
\begin{tabular}{lccccc}
\toprule
Method & PSNR$\uparrow$ & SSIM$\uparrow$ & LPIPS$\downarrow$ & Depth L1 (m)$\downarrow$ \\
\midrule
\StreamSplat~\cite{wu2026streamsplat} & \TODO{} & \TODO{} & \TODO{} & --- \\
\StereoSplat (no metric) & \TODO{} & \TODO{} & \TODO{} & --- \\
\StereoSplat (ours) & \TODO{} & \TODO{} & \TODO{} & \TODO{} \\
\StereoSplat (oracle) & \TODO{} & \TODO{} & \TODO{} & \TODO{} \\
\bottomrule
\end{tabular}
\end{table}

\subsection{Ablation Study (Q2)}

We ablate five axes (Table~\ref{tab:ablation}).
\textbf{(A)} Depth source: stereo disparity vs.\ relative depth vs.\ GT.
\textbf{(B)} Input views: $V{=}2$ (stereo) vs.\ $V{=}1$ (mono).
\textbf{(C)} Rasterizer: perspective vs.\ orthographic.
\textbf{(D)} Display operators: smoothing on/off, popping suppression
on/off.
\textbf{(E)} Loss terms: with/without $\Ldepth$, with/without $\Lstereo$.

\begin{table}[t]
\centering
\small
\caption{Ablation. Each row toggles one axis relative to the full model.
\TODO{numbers to be filled.}}
\label{tab:ablation}
\begin{tabular}{lcccc}
\toprule
Configuration & PSNR$\uparrow$ & Depth L1$\downarrow$ & Temp. std$\downarrow$ \\
\midrule
Full model & \TODO{} & \TODO{} & \TODO{} \\
A2: relative depth & \TODO{} & --- & \TODO{} \\
B2: mono input ($V{=}1$) & \TODO{} & \TODO{} & \TODO{} \\
C2: orthographic rasterizer & \TODO{} & --- & \TODO{} \\
D4: no display operators & \TODO{} & \TODO{} & \TODO{} \\
E2: no $\Ldepth$ & \TODO{} & \TODO{} & \TODO{} \\
E3: no $\Lstereo$ & \TODO{} & \TODO{} & \TODO{} \\
\bottomrule
\end{tabular}
\end{table}

The ablation confirms that the metric depth source is the sole source of
metric scale (A2 reports no depth error), the stereo input is necessary
for the right-view consistency loss to function (B2), the perspective
rasterizer is necessary for metric poses to render correctly (C2), and
both $\Ldepth$ and $\Lstereo$ contribute to the final quality (E2, E3).

\subsection{Display Operator Study (Q3)}

Figure~\ref{fig:display} shows spatiotemporal images (a vertical slice
over time, following~\cite{yun2025compensating}) for the four display
configurations: raw, smoothing only, popping suppression only, and both.

\begin{figure}[t]
\centering
\framebox[0.46\textwidth]{\parbox{0.44\textwidth}{\centering\vspace{2pt}
Spatiotemporal-image comparison across the four display configurations.
\vspace{2pt}}}
\caption{Display-operator study. Static-region flicker is visibly reduced
when both operators are active.}
\label{fig:display}
\end{figure}

Without any operator, static regions exhibit visible flicker; with both
operators, the spatiotemporal image is markedly cleaner. Quantitatively,
the temporal standard deviation in static regions drops from
\TODO{X} (raw) to \TODO{X} (both), with negligible change in PSNR,
confirming that the operators improve temporal consistency without
harming reconstruction quality.

\subsection{Qualitative Results (Q4)}

Figure~\ref{fig:qual} shows side-by-side renderings from
\StreamSplat (monocular, orthographic) and \StereoSplat on a DSEC driving
sequence.

\begin{figure}[t]
\centering
\framebox[0.46\textwidth]{\parbox{0.44\textwidth}{\centering\vspace{2pt}
Qualitative side-by-side: \StreamSplat vs.\ \StereoSplat on a DSEC
sequence, with novel-view re-rendering.
\vspace{2pt}}}
\caption{Qualitative comparison. \StereoSplat re-renders correctly under
viewpoint shifts thanks to the perspective rasterizer.}
\label{fig:qual}
\end{figure}

\StereoSplat produces a metric scene that can be re-rendered at
arbitrary novel viewpoints with correct perspective, while \StreamSplat's
orthographic canonical space produces visible distortions under viewpoint
shift. Figure~\ref{fig:depth} compares the rendered metric depth against
the IGEV-Stereo depth, showing sub-meter agreement on most of the image.

\begin{figure}[t]
\centering
\framebox[0.46\textwidth]{\parbox{0.44\textwidth}{\centering\vspace{2pt}
Rendered metric depth (left) vs.\ IGEV-Stereo depth (right).
\vspace{2pt}}}
\caption{Depth comparison on DSEC.}
\label{fig:depth}
\end{figure}

Please see the supplementary video for animated comparisons.

\paragraph{Implementation.}
All experiments run on a single NVIDIA RTX 4090. Training uses AdamW~\cite{loshchilov2019decoupled} with
learning rate $5{\times}10^{-4}$, batch size 8, and the two-stage schedule
of \StreamSplat. The display module adds less than 1\,ms per frame and is
training-free. Code, configs, and pretrained checkpoints will be released.
